                                     IMC2011, Blagoevgrad, Bulgaria
                                                  Day 1, July 30, 2011



Problem 1. Let f : R → R be a continuous function. A point x is called a shadow point if there exists a point y ∈ R with
y > x such that f (y) > f (x). Let a < b be real numbers and suppose that
   • all the points of the open interval I = (a, b) are shadow points;
   • a and b are not shadow points.
Prove that
   a) f (x) ≤ f (b) for all a < x < b;
   b) f (a) = f (b).
                                                                                            (José Luis Dı́az-Barrero, Barcelona)
Solution. (a) We prove by contradiction. Suppose that exists a point c ∈ (a, b) such that f (c) > f (b).
    By Weierstrass’ theorem, f has a maximal value m on [c, b]; this value is attained at some point d ∈ [c, b]. Since
f (d) = max f ≥ f (c) > f (b), we have d 6= b, so d ∈ [c, b) ⊂ (a, b). The point d, lying in (a, b), is a shadow point, therefore
        [c.b]
f (y) > f (d) for some y > d. From combining our inequalities we get f (y) > f (d) > f (b).
    Case 1: y > b. Then f (y) > f (b) contradicts the assumption that b is not a shadow point.
    Case 2: y ≤ b. Then y ∈ (d, b] ⊂ [c, b], therefore f (y) > f (d) = m = max f ≥ f (y), contradiction again.
                                                                             [c,b]

   (b) Since a < b and a is not a shadow point, we have f (a) ≥ f (b).
   By part (a), we already have f (x) ≤ f (b) for all x ∈ (a, b). By the continuity at a we have

                                            f (a) = lim f (x) ≤ lim f (b) = f (b)
                                                    x→a+0           x→a+0

   Hence we have both f (a) ≥ f (b) and f (a) ≤ f (b), so f (a) = f (b).

Problem 2. Does there exist a real 3 × 3 matrix A such that tr(A) = 0 and A2 + At = I? (tr(A) denotes the trace of A,
At is the transpose of A, and I is the identity matrix.)
                                                                                         (Moubinool Omarjee, Paris)
Solution. The answer is NO.
   Suppose that tr(A) = 0 and A2 + At = I. Taking the transpose, we have

                                    A = I − (A2 )t = I − (At )2 = I − (I − A2 )2 = 2A2 − A4 ,

                                                      A4 − 2A2 + A = 0.
                                                                                     √
The roots of the polynomial x4 − 2x2 + x = x(x − 1)(x2 + x − 1) are 0, 1, −1±  2
                                                                                 5
                                                                                   so these numbers can be the eigenvalues of
                                         √
                                      1±   5
A; the eigenvalues of A2 can be 0, 1, 2 .
    By tr(A) = 0, the sum of the eigenvalues is 0, and by tr(A2 ) = tr(I − At ) = 3 the sum of squares of the eigenvalues is 3.
It is easy to check that this two conditions cannot be satisfied simultaneously.

Problem 3. Let p be a prime number. Call a positive integer n interesting if

                                               xn − 1 = (xp − x + 1)f (x) + pg(x)

for some polynomials f and g with integer coefficients.
    a) Prove that the number pp − 1 is interesting.
    b) For which p is pp − 1 the minimal interesting number?
                                                                           (Eugene Goryachko and Fedor Petrov, St. Petersburg)
Solution. (a) Let’s reformulate the property of being interesting: n is interesting if xn − 1 is divisible by xp − x + 1 in the
ring of polynomials over Fp (the field of residues modulo p). All further congruences are modulo xp − x + 1 in this ring. We
                          2                                         3      2
have xp ≡ x − 1, then xp = (xp )p ≡ (x − 1)p ≡ xp − 1 ≡ x − 2, xp = (xp )p ≡ (x − 2)p ≡ xp − 2p ≡ x − 2p − 1 ≡ x − 3 and
                                             p
so on by Fermat’s little theorem, finally xp ≡ x − p ≡ x,
                                                            p
                                                       x(xp −1 − 1) ≡ 0.
                                                                       p
Since the polynomials xp − x + 1 and x are coprime, this implies xp −1 − 1 ≡ 0.


                                                                1
   (b) We write
                          2        p−1                   2         p−1
                    x1+p+p +···+p        = x · xp · xp · . . . · xp      ≡ x(x − 1)(x − 2) . . . (x − (p − 1)) = xp − x ≡ −1,
                2        p−1
hence x2(1+p+p +···+p ) ≡ 1 and a = 2(1 + p + p2 + · · · + pp−1 ) is an interesting number.
                         2
    If p > 3, then a = p−1 (pp − 1) < pp − 1, so we have an interesting number less than pp − 1. On the other hand, we show
that p = 2 and p = 3 do satisfy the condition. First notice that by gcd(xm − 1, xk − 1) = xgcd(m,k) − 1, for every fixed p the
greatest common divisors of interesting numbers is also an interesting number. Therefore the minimal interesting number
divides all interesting numbers. In particular, the minimal interesting number is a divisor of pp − 1.
    For p = 2 we have pp − 1 = 3, so the minimal interesting number is 1 or 3. But x2 − x + 1 does not divide x − 1, so 1 is
not interesting. Then the minimal interesting number is 3.
    For p = 3 we have pp − 1 = 26 whose divisors are 1, 2, 13, 26. The numbers 1 and 2 are too small and x13 ≡ −1 6≡ +1 as
shown above, so none of 1,2 and 13 is interesting. So 26 is the minimal interesting number.
    Hence, pp − 1 is the minimal interesting number if and only if p = 2 or p = 3.

Problem 4. Let A1 , A2 , . . . , An be finite, nonempty sets. Define the function
                                                      n
                                                      X            X
                                            f (t) =                              (−1)k−1 t|Ai1 ∪Ai2 ∪...∪Aik | .
                                                      k=1 1≤i1 <i2 <...<ik ≤n

Prove that f is nondecreasing on [0, 1].
   (|A| denotes the number of elements in A.)
                                                                                           (Levon Nurbekyan and Vardan Voskanyan, Yerevan)
                           n
                           S
Solution 1. Let Ω =              Ai . Consider a random subset X of Ω which chosen in the following way: for each x ∈ Ω, choose
                           i=1
the element x for the set X with probability t, independently from the other elements.
   Then for any set C ⊂ Ω, we have
                                                     P (C ⊂ X) = t|C| .
By the inclusion-exclusion principle,
                                                                                         
                                              P (A1 ⊂ X) or (A2 ⊂ X) or . . . or (An ⊂ X) =
                                         n
                                         X             X
                                                                 (−1)k−1 P Ai1 ∪ Ai2 ∪ . . . ∪ Aik ⊂ X =
                                                                                                      
                                    =
                                         k=1 1≤i1 <i2 <...<ik ≤n
                                                   n
                                                   X           X
                                               =                           (−1)k−1 t|Ai1 ∪Ai2 ∪...∪Aik | .
                                                   k=1 1≤i1 <i2 <...<ik ≤n
                                               
The probability P (A1 ⊂ X) or . . . or (An ⊂ X) is a nondecreasing function of the probability t.

Problem 5. Let n be a positive integer and let V be a (2n − 1)-dimensional vector space over the two-element field.
Prove that for arbitrary vectors v1 , . . . , v4n−1 ∈ V , there exists a sequence 1 ≤ i1 < . . . < i2n ≤ 4n − 1 of indices such that
vi1 + . . . + vi2n = 0.
                                                                               (Ilya Bogdanov, Moscow and Géza Kós, Budapest)
Solution. Let V = aff{v1 , . . . , v4n−1 }. The statement vi1 + · · · + vi2n = 0 is translation-invariant (i.e. replacing the vectors
by v1 − a, . . . , v4n−1 − a), so we may assume that 0 ∈ V . Let d = dim V .
Lemma. The vectors can be permuted in such a way that v1 + v2 , v3 + v4 , . . . , v2d−1 + v2d form a basis of V .
Proof. We prove by induction on d. If d = 0 or d = 1 then the statement is trivial.
   First choose the vector v1 such a way that aff(v2 , v3 , . . . , v4n−1 ) = V ; this is possible since V is generated by some d + 1
vectors and we have d + 1 ≤ 2n < 4n − 1. Next, choose v2 such that v2 6= v1 . (By d > 0, not all vectors are the same.)
   Now let ℓ = {0, v1 + v2 } and let V ′ = V /ℓ. For any w ∈ V , let w̃ = ℓ + w = {w, w + v1 + v2 } be the class of the factor
space V ′ containing w. Apply the induction hypothesis to the vectors v˜3 , . . . , v4n−1 ˜ . Since dim V ′ = d − 1, the vectors can
                                                                            ′
permuted in such a way that v˜3 + v˜4 , . . . , v2d−1
                                                  ˜ + v˜2d is a basis of V . Then v1 + v2 , v3 + v4 , . . . , v2d−1 + v2d is a basis of V .
      Now we can assume that v1 + v2 , v3 + v4 , . . . , v2d−1 + v2d is a basis of V . The vector w = (v1 + v3 + · · · + v2d−1 ) + (v2d+1 +
v2d+2 + · · · + v2n+d ) is the sum of 2n vectors, so w ∈ V . Hence, w + ε1 (v1 + v2 ) + · · · + εd (v2d−1 + v2d ) = 0 with some
ε1 , . . . , εd ∈ F2 , therefore
                                           X d                                 2n+d
                                                                                  X
                                                    (1 − εi )v2i−1 + εi v2i +           vi = 0.
                                                   i=1                                   i=2d+1

The left-hand side is the sum of 2n vectors.

                                                                             2
